%!TEX root = ./../../thesis-trier.tex

\IN{solvable automaton}
We call $\mathcal A$ \emph{solvable}
if $\delta(q, ab) \lereach q$ implies $\delta(q, ab) = \delta(q, ba)$
for all $q \in Q$ and $a,b \in \Sigma$.
Clearly, every commutative automaton is solvable.
Since for partially ordered automata,
$\delta(q, ab) \lereach q$
implies $\delta(q, a) = \delta(q, b) = q$,
and so $\delta(q, ab) = \delta(q, ba) = q$,
partially ordered automata are also solvable. Clearly, commutative automata
are solvable as well.

Here, we investigate solvable automata. In \cref{prelim:classes:sec:solvable:operations}, 
we show that the class of solvable automata
is closed under various operations. \Cref{prelim:classes:sec:solvable:char}
gives other characterizations of solvability.
In \cref{prelim:classes:sec:solvable:sync}, 
we prove bounds on shortest words whose rank is bounded,
and we give a technical result on words mapping one subset into another subset.
Solvable automata were introduced by
Rystsov~\cite{Rystsov97} with a slightly different, more restricted,
definition, we discuss this in \cref{prelim:classes:sec:comparison}.

\subsection{Operations}
\label{prelim:classes:sec:solvable:operations}

Here, we prove that solvability
is preserved 
in subautomata and morphic images, Rees quotient, product, and pair automata
in \cref{prelim:classes:solvable:prop:subaut},
\cref{prelim:classes:solvable:prop:morph}, 
\cref{prelim:classes:solvable:cor:subhom},
\cref{prelim:subclasses:prop:solvable-product}, and
\cref{prelim:subclasses:cor:solv-pair}, respectively.
\Cref{prelim:subclasses:cor:solvable-product}
proves that 
solvability of $\mathcal A$, $\mathcal A \times \mathcal A$, and $\mathcal A^{\le 2}$ are equivalent properties,
and in \cref{prelim:subclasses:cor:solv-pair} 
that solvability of $\mathcal A$
implies solvability of 
 $\mathcal A^{\le 2} / {\sim_F}$
and 
$( \mathcal A \times \mathcal A ) / {\sim_D}$
with $F = \{ \{q\} \mid q \in Q \}$
and $D = \{ (q,q) \mid q \in Q \}$.
In \cref{prelim:classes:solv:prop:perm-product-solv-implies-solv}
we show that for permutation automata,
solvability of~$( \mathcal A \times \mathcal A ) / {\sim_D}$
implies it for~$\mathcal A$.

\begin{proposition} \label{prelim:classes:solvable:prop:subaut}
	Let $\semiautomat{A}{Q}{\Sigma}{\delta}$ be a solvable automaton
	and $S \subseteq Q$ be an upper set.
Then $\mathcal A_{|S}$ is solvable.	
\end{proposition}
\begin{proof}
	Suppose $\delta_{|S}(q, ab) \lereach q$
	for some $q \in S$ and $a,b \in \Sigma$.
	Then $\delta(q, ab) \lereach q$
	and, by solvability of $\mathcal A$,
	we have $\delta_{|S}(q, ab) = \delta(q, ab) = \delta(q, ba) = \delta_{|S}(q, ba)$.
	Hence, $\mathcal A_{|S}$ is solvable.
\end{proof} 

\begin{proposition} \label{prelim:classes:solvable:prop:morph}
	Let $\semiautomat{A}{Q}{\Sigma}{\delta}$
	and $\semiautomat{B}{S}{\Sigma}{\mu}$.
	Suppose there exists a surjective morphism $\varphi \colon Q \to S$.
	Then solvability of $\mathcal A$ implies solvability of~$\mathcal B$.
\end{proposition}
\begin{proof}
	Suppose we have $\mu(\varphi(q), abu) = \varphi(q)$ for some $u \in \Sigma^*$.
	Then it is $\varphi(\delta(q, (abu)^i)) = \varphi(q)$
	for all $i \ge 0$. By finiteness, there exist $i \ge 0$
	and $p > 0$
	such that $\delta(q, (abu)^{i + p}) = \delta(q, (abu)^i)$.
	Hence $\delta(q, (abu)^{i}ab) \lereach \delta(q, (abu)^i)$.
	Since~$\mathcal A$ is solvable,
	we find $\delta(q, (abu)^iab) = \delta(q, (abu)^iba)$.
	Then 
	\begin{multline*}
		\mu(\varphi(q), ab)  = \mu(\varphi(\delta(q, (abu)^i)), ab) 
		= \varphi(\delta(q, (abu)^i ab)) \\
		= \varphi(\delta(q, (abu)^i ba)) 
		= \mu(\varphi(\delta(q, (abu)^i)),ba) 
		= \mu(\varphi(q), ba).
	\end{multline*}	
	Hence, $\mathcal B$ is solvable.
\end{proof}

By \cref{prelim:classes:solvable:prop:morph} the class of solvable automata %, as well as all the other classes of automata introduced, 
is closed under isomorphisms.

\begin{corollary}\label{prelim:classes:solvable:cor:subhom}
	Let $\semiautomat{A}{Q}{\Sigma}{\delta}$ be a solvable automaton
	and $S \subseteq Q$ be an upper set.
	Then $\mathcal A / {\sim_S}$ is solvable.
\end{corollary} 
\begin{proof}
	The canonical morphism for $\sim_S$
	from \cref{prelim:prop:congruence_induces_morphism} is surjective and
	maps $\mathcal A$ onto $\mathcal A / {\sim_S}$.
	So the claim follows from \cref{prelim:classes:solvable:prop:morph}.
\end{proof}

\begin{proposition}\label{prelim:subclasses:prop:solvable-product}
	Let $\semiautomat{A_1}{Q_1}{\Sigma}{\delta_1}$ and 
	$\semiautomat{A_2}{Q_2}{\Sigma}{\delta_2}$.
	Then $\mathcal A_1$ and $\mathcal A_2$ are solvable
	if and only if $\mathcal A_1 \times \mathcal A_2$ is solvable.
\end{proposition}
\begin{proof}
	First, assume $\mathcal A_1 \times \mathcal A_2$ is solvable.
	Since the maps $(p,q) \mapsto p$
	and $(p,q) \mapsto q$
	are surjective morphisms from $\mathcal A_1 \times \mathcal A_2$
	to $\mathcal A_1$ and~$\mathcal A_2$, respectively,
	we find that $\mathcal A_1$
	and~$\mathcal A_2$ are solvable by \cref{prelim:classes:solvable:prop:morph}. 
	
	Now, assume $\mathcal A_1$ and $\mathcal A_2$
	are solvable. 
	Let $(p,q) \in Q_1 \times Q_2$ and $a,b \in \Sigma$
	be such that $\delta_{\mathcal A_1 \times \mathcal A_2}((p,q), ab) \lereach (p,q)$.
	This gives $\delta_1(p, ab) \lereach p$
	and $\delta_2(q, ab) \lereach q$.\footnote{Note that
		$\delta_1(p, ab) \lereach p$
		and $\delta_2(q, ab) \lereach q$
		not necessarily imply $\delta_{\mathcal A_1 \times \mathcal A_2}((p,q), ab) \lereach (p,q)$.}
	Since $\mathcal A_1$ and $\mathcal A_2$ are solvable,
	we find $\delta_1(q, ab) = \delta_1(q, ba)$
	and $\delta_2(p, ab) = \delta_2(q, ba)$.
	Thus, 
	we get $\delta_{\mathcal A_1 \times \mathcal A_2}((p,q), ab) = \delta_{\mathcal A_1 \times \mathcal A_2}((p,q), ba)$
	from the definition of~$\mathcal A_1 \times \mathcal A_2$.
\end{proof}

\begin{corollary}\label{prelim:subclasses:cor:solvable-product}
	Let $\semiautomat{A}{Q}{\Sigma}{\delta}$ and $i \ge 2$.
	Then the following are equivalent.
	
	\begin{enumerate}
		\item \label{solvable-product:item:1} $\mathcal A$ is solvable.
		
		\item \label{solvable-product:item:2a}
		$\mathcal A \times \cdots \times \mathcal A$ ($i$ times) is solvable.
		
		\item \label{solvable-product:item:2} $\mathcal A \times \mathcal A$ is solvable.
		
		\item \label{solvable-product:item:3} $\mathcal A^{\le 2}$ is solvable. % als subautomton einbetten
	\end{enumerate}	
\end{corollary}
\begin{proof}
	$\ref{solvable-product:item:1} \Rightarrow \ref{solvable-product:item:2a}$: This follows from \cref{prelim:subclasses:prop:solvable-product}
	applied $i$ times.
	
	$\ref{solvable-product:item:2a} \Rightarrow \ref{solvable-product:item:2}$:
	We can assume $i > 2$.
	Then 
	\[
	\mathcal A \times \cdots \times \mathcal A = \underbrace{(\mathcal A \times \cdots \times \mathcal A)}_{\text{$i-2$ times}} \times (\mathcal A \times \mathcal A).
	\]
	By \cref{prelim:subclasses:prop:solvable-product} 
	we obtain that $\mathcal A \times \mathcal A$
	is solvable.
			
	$\ref{solvable-product:item:2} \Rightarrow \ref{solvable-product:item:3}$: We have $\mathcal A^{\le 2} \cong (\mathcal A \times \mathcal A) / \sim$
	for the $\TS{\mathcal A \times \mathcal A}$-congruence $\sim$
	given by $(p_1,p_2) \sim (q_1, q_2)$
	if and only if $\{p_1, p_2\} = \{q_1, q_2\}$.
	%with $T = \{ ((p,q), (q,p)) \mid p,q \in Q\}$.
	So, the claim follows from \cref{prelim:classes:solvable:cor:subhom}. 
		
	$\ref{solvable-product:item:3} \Rightarrow \ref{solvable-product:item:1}$: % Follows from \cref{prelim:classes:solvable:prop:subhom}, 
	As $(\mathcal A^{\le 2})_{|F} \cong \mathcal A$
	with $F = \{ \{q\} \mid q \in F\}$, the claim follows from \cref{prelim:classes:solvable:cor:subhom}.
\end{proof} 

\begin{corollary} \label{prelim:subclasses:cor:solv-pair}
	Let $\semiautomat{A}{Q}{\Sigma}{\delta}$ be solvable.
	Then the automata $\mathcal A^{\le 2} / {\sim_F}$
	and 
	$( \mathcal A \times \mathcal A ) / {\sim_D}$
	with $F = \{ \{q\} \mid q \in Q \}$
	and $D = \{ (q,q) \mid q \in Q \}$
	are solvable.	
\end{corollary}
\begin{proof}
	For solvable $\semiautomat{A}{Q}{\Sigma}{\delta}$,
	$\mathcal A^{\le 2}$ and $\mathcal A \times \mathcal A$ are solvable by \cref{prelim:subclasses:cor:solvable-product}.
	So, \cref{prelim:classes:solvable:cor:subhom}
    gives the claim.
\end{proof}

The converse statement to \cref{prelim:subclasses:cor:solv-pair}
does not hold in general, as the example from \cref{prelim:classes:fig:solvablenotR2solvable}
shows.

\begin{figure}[ht]
	\centering
	\begin{subfigure}{0.33\textwidth}
		\centering
		\begin{tikzpicture}[node distance = 2cm, on grid, auto, ->]
			\node (q1) [state] {$q_1$};
			\node (q2) [state, right of = q1] {$q_2$};
			
			%	\path[->] 
			%	(q2) edge [loop above] node {$a$} (q2);
			\path[->] % (q1) [loop above] edge node {$b$} (q1)
			(q1) edge [bend left] node {$a, b$} (q2)
			(q2) edge [bend left, below] node {$b$} (q1); 
		\end{tikzpicture}
		\caption{The automaton $\mathcal A$.}
	\end{subfigure}%
	\begin{subfigure}{0.4\textwidth}
		\centering
		\begin{tikzpicture}[node distance = 2.2cm, on grid, auto, ->]
			\tikzstyle{state-subset}=[]
			
			\node (12) [state-subset] {$(q_1, q_2)$};
			\node (D)  [state-subset, right of = 12] {$D$};
			\node (21) [state-subset, right of = D] {$(q_2, q_1)$};
			
			\path[->] (12) edge [bend left, out=48,in=132] node {$b$} (21)
			(21) edge [bend left, out=48,in=132] node {$b$} (12);
			\path[->] (12) edge node {$a$} (D)
			(21) edge node {$a$} (D);
			% (D)  edge [loop above] node {$a,b$} (D);
		\end{tikzpicture}
		\caption{$( \mathcal A \times \mathcal A ) / {\sim_D}$, which is solvable.}
	\end{subfigure}%
	\begin{subfigure}{0.27\textwidth}
		\centering
		\begin{tikzpicture}[node distance = 1.5cm, on grid, auto, ->]
			\tikzstyle{state-subset}=[]
			
			\node (12) [state-subset] {$\{q_1, q_2\}$};
			\node (D)  [state-subset, below of = 12] {$F$};
			
			\path[->] (12) edge node {$a$} (D);
			\path[->] (12) edge [loop above]  node {$b$} (12);
		\end{tikzpicture}
		\caption{$\mathcal A^{\le 2} / {\sim_F}$}
	\end{subfigure}%
	\caption[Non-solvable automaton with solvable pair automaton]{The automaton $\mathcal A$ is not solvable, but $( \mathcal A \times \mathcal A ) / {\sim_D}$ and $\mathcal A^{\le 2} / {\sim_F}$ with $D = \{(q,q) \mid q \in Q\}$ and $F = \{\{q\}\mid q \in Q\}$ are solvable. Self-loops omitted.}
	\label{prelim:classes:fig:solvablenotR2solvable}
\end{figure} 

However, for permutation automata,
a converse of \cref{prelim:subclasses:cor:solv-pair} holds
for $( \mathcal A \times \mathcal A ) / {\sim_D}$.

\begin{proposition} \label{prelim:classes:solv:prop:perm-product-solv-implies-solv}
	Let $\semiautomat{A}{Q}{\Sigma}{\delta}$ be a permutation automaton.
	Then solvability of $( \mathcal A \times \mathcal A ) / {\sim_D}$ with $D = \{ (q,q) \mid q \in Q \}$
	implies that $\mathcal A$ is solvable.	
\end{proposition}
\begin{proof}
	We can assume $|Q| \ge 2$.
	Clearly, $\mathcal A$ is solvable
	if each subautomaton of the form $\mathcal A_{|\upperset{\{q\}}}$
	for each $q \in Q$ is solvable.
	Thus, we can assume $Q = \upperset{\{q\}}$ for some $q \in Q$
	in the following.
	Let $p \in Q$ be a state with $p \ne q$.
	Consider the upper set $S = \upperset{\{(p, q)\}}$
	in $\mathcal A \times \mathcal A$.
	Since $p \ne q$ and $\mathcal A$ is a permutation automaton, we have
	\[
	( \mathcal A \times \mathcal A )_{|S} \cong \left( (\mathcal A \times \mathcal A) / {\sim_D} \right)_{|\pi_{\sim_D}(S)}.
	\]
	and the projection mapping restricted to~$S$, i.e., $(\pi_2)_{|S} S \to Q$,
	is a surjective morphism from $ ( \mathcal A \times \mathcal A )_{|S}$
	onto $\mathcal A$.
	Thus, by \cref{prelim:classes:solvable:prop:subaut,prelim:classes:solvable:prop:morph} 
	we find that~$\mathcal A$ is solvable.	
\end{proof}

The partition $\{ \{ (p,q), (q, p) \} \} \cup \{ \{ D \} \}$
of the state set of $(\mathcal A \times \mathcal A ) / {\sim_D}$
induces a congruence whose quotient automaton
is isomorphic to $\mathcal A^{\le 2} / {\sim_F}$. 
Then by \cref{prelim:classes:solvable:prop:morph,prelim:prop:congruence_induces_morphism}
it follows that $\mathcal A^{\le 2} / {\sim_F}$
is solvable if $(\mathcal A \times \mathcal A ) / {\sim_D}$ is solvable.
However, we do not know whether the converse implication holds.

\subsection{Characterizations}
\label{prelim:classes:sec:solvable:char}

Solvable automata enjoy a wealth of other characterizations
as we show next. \Cref{prelim:classes:solvable:thm} 
collects various characterizations in terms of strongly
connected components. \Cref{prelim:classes:solv:prop:char-seq} 
characterizes solvability by increasing sequences of upper
sets, resembling one definition originally given by Rystsov~\cite{Rystsov97} 
(see the discussion in \cref{prelim:classes:sec:comparison}).
\Cref{prelim:classes:solvable:prop:char-inductive} 
gives definition in terms of a single upper sets,
some of them are suitable for an inductive reasoning and construction
of solvable automata, starting from commutative automata.

\begin{theorem}\label{prelim:classes:solvable:thm}
Let $\semiautomat{A}{Q}{\Sigma}{\delta}$.
Then the following are equivalent.
\begin{enumerate}
\item \label{prelim:classes:solvable:thm:char:solvable} $\mathcal A$ is solvable.

\item \label{prelim:classes:solvable:thm:char:SCC-quotients-comm}  For each $q \in Q$, the automaton
\[
( \mathcal A_{|S} ) / {\sim_T}
\]
with $S = \upperset{\{q\}}$ and $T = S \setminus [q]_{\eqreach}$ % mit isomorphiesätzen dann auf reihen
is commutative.

\item \label{prelim:classes:solvable:thm:char:SCC-quotients-comm-perm}  For each $q \in Q$, the automaton
\[ 
( \mathcal A_{|S} ) / {\sim_T}
\]
with $S = \upperset{\{q\}}$ and $T = S \setminus [q]_{\eqreach}$ % mit isomorphiesätzen dann auf reihen
is a strongly trap-connected commutative permutation-reset automaton
or a commutative permutation automaton (where the latter case includes~$T = \emptyset$, in which case~$\mathcal A$
is strongly connected).

\item \label{prelim:subclasses:thm:solvable:permutation-reset} For all upper sets $S, T \subseteq Q$
such that $T \subseteq S$ and $S \setminus T$ 
is a strongly connected component, the automaton $( \mathcal A_{|S} ) / {\sim_T}$
is a commutative permutation-reset automaton.




\item \label{prelim:subclasses:thm:solvable:restrict} For each strongly connected component $S \subseteq Q$
the automaton $\restrictAutGamS{\mathcal A}{\Gamma_1}{S}$
with $\Gamma_1 = \{ a \in \Sigma \mid \delta(S, a) = S \}$
is a commutative permutation automaton 
and $\Sigma = \Gamma_1 \cup \Gamma_2$
for $\Gamma_2 = \{ a \in \Sigma \mid \delta(S, a) \cap S = \emptyset \}$.



\item \label{prelim:prop:solvable-SCC:2} Each strongly connected component $S \subseteq Q$ is a block
such that $\delta(q, ab) = \delta(q, ba)$
for all $q \in S$ and $a,b \in \Sigma$
with $\delta(S, a) = \delta(S, b) = S$.

\item \label{prelim:prop:solvable-SCC:3a} For each strongly connected component $S \subseteq Q$
and $a,b \in \Sigma$ with $\delta(S, a) \cap S \ne \emptyset$
and $\delta(S, b) \cap S \ne \emptyset$
we have $\delta(q, ab) = \delta(q, ba)$ for all $q \in S$.

\item \label{prelim:prop:solvable-SCC:3} For each strongly connected $S \subseteq Q$
and $a,b \in \Sigma$ with $\delta(S, a) \cap S \ne \emptyset$
and $\delta(S, b) \cap S \ne \emptyset$ we have $\delta(q, ab) = \delta(q, ba)$ for all $q \in S$.

\item \label{prelim:prop:solvable-SCC:4} For each strongly connected $S \subseteq Q$
we have $\delta(q, ab) = \delta(q, ba)$
for all $q \in S$ and all $a,b \in \Sigma$ with $\delta(q, ab) \in S$.

\item \label{prelim:prop:solvable-SCC:3b} For each $q \in Q$ and $u \in \Sigma^*$, $\delta(q,u) = q$
 implies $\delta(q, v) = \delta(q, u)$ for all $v \in \Sigma^*$
 such that $|v|_x = |u|_x$ for all $x \in \Sigma$. 
 
\item \label{prelim:prop:solvable-SCC:3c} For each $q \in Q$ and $u \in \Sigma^*$, $\delta(q,u) = q$
implies
\[
\delta(q, v) = \delta(q, v_{\pi(1)} \cdots v_{\pi(m)})
\]
for each permutation $\pi\colon \{1,\ldots,m\} \to \{1,\ldots,m\}$ 
and $v \in \Sigma^*$ with $u \in v\Sigma^*$
and $v = v_1\cdots v_m$ for $v_1, \ldots, v_m \in \Sigma$.
\end{enumerate}
\end{theorem}
\begin{proof}
$\ref{prelim:classes:solvable:thm:char:solvable} \Rightarrow \ref{prelim:classes:solvable:thm:char:SCC-quotients-comm}$: 	
Let $p \in S$ with $p \notin T$. Assume $\delta_{|S}(p, ab) \ne \delta_{|S}(p, ab)$.
Then $\delta(p, ab) \ne \delta(p, ba)$
and since~$\mathcal A$
is solvable, we have $\delta(p, ab) \not\eqreach p$ and $\delta(p, ab) \not\eqreach p$.
Since $p \eqreach q$, we find $\delta(p, ab) \not\eqreach q$ 
and  $\delta(p, ba) \not\eqreach q$ 
and so $\delta(p, ab) , \delta(p, ba)\in T$ by \cref{prelim:prop:reachability-finest,prelim:eqn:reach-upperset}.
Hence, in $(\mathcal A_{|S} ) / {\sim_T}$,
the state~$p$ is mapped to~$T$ by both words $ab$ and $ba$.
This already establishes commutativity
of $(\mathcal A_{|S} ) / {\sim_T}$
since if  $\delta_{|S}(p, ab) = \delta_{|S}(p, ab)$,
then also in $(\mathcal A_{|S} ) / {\sim_T}$
the state~$p$ is mapped to the same state
by~$ab$ and~$ba$.

$\ref{prelim:classes:solvable:thm:char:SCC-quotients-comm} \Rightarrow \ref{prelim:classes:solvable:thm:char:SCC-quotients-comm-perm}$: 	
By assumption~$(\mathcal A_{|S} ) / {\sim_T}$
is commutative. Hence, each letter acts as a morphism 
on its state set, which is $[q]_{\eqreach} \cup \{T\}$.
By \cref{prelim:automata:cor:SCC-disjoint-or-permutates},
each letter either permutes $[q]_{\eqreach}$ 
or maps all states in it to $T$. So,
$(\mathcal A_{|S} ) / {\sim_T}$
is a permutation-reset automaton.
If there exists a letter mapping $[q]_{\simeq}$
to~$T$, then it is not a sink component
and by \crefitem{prelim:automata:prop:char-strongly-trap}{prelim:automata:prop:char-strongly-trap:1b}
it is strongly trap-connected.
Otherwise, since~$T$ is a trap state for $T \ne \emptyset$,
each letter permutes all states.

Alternatively, this follows with \cref{cor:SCC_comm_permutation_automaton}.

$\ref{prelim:classes:solvable:thm:char:SCC-quotients-comm-perm} \Rightarrow \ref{prelim:subclasses:thm:solvable:permutation-reset}$:
First, assume~$S$ is a sink component.
Then, by Item~\ref{prelim:classes:solvable:thm:char:SCC-quotients-comm-perm},
$\mathcal A_{|S}$ is a commutative permutation automaton.
So, $\mathcal (\mathcal A_{|S}) / {\sim_T}$
is a commutative permutation automaton, since for $T \ne \emptyset$ 
merely a trap state is added.
Next, assume~$S$ is not a sink component. 
By assumption, we have $S \setminus T = [q]_{\eqreach}$
for some $q \in Q$. Since~$S$ is an upper set, we have 
$\upperset{\{q\}} \subseteq S \subseteq T$. 
Set $K = \upperset{\{q\}}$ and $R = K \setminus [q]_{\eqreach}$,
where $R \ne \emptyset$ and $T \ne \emptyset$
as~$S$ is a sink component.
Then, it is
\[
 ( \mathcal A_{|S} ) / {\sim_T}
 \cong 
 ( \mathcal A_{|K} ) / {\sim_{R}}
\]
since $\sim_T$ on~$S$
equals $[q]_{\eqreach} \times [q]_{\eqreach} \cup \{ T \}$
and $\sim_R$ on~$K$ equals $[q]_{\eqreach} \times [q]_{\eqreach} \cup \{ R \}$
and so both induce isomorphic Rees quotient automata.
By Item~\ref{prelim:classes:solvable:thm:char:SCC-quotients-comm-perm},
$( \mathcal A_{|K} ) / {\sim_{R}}$
is a commutative permutation-reset automaton.

$\ref{prelim:subclasses:thm:solvable:permutation-reset} \Rightarrow \ref{prelim:subclasses:thm:solvable:restrict}$:
Let $S \subseteq Q$ be a strongly connected component. By assumption,
$( \mathcal A_{|R} ) / {\sim_T}$ with $R = \upperset{S}$ and $T = R \setminus S$
is a permutation-reset automaton, which is equivalent
to~$S$ being a block by \cref{prelim:perm-reset-Rees-factor}, 
which gives $\Sigma = \Gamma_1 \cup \Gamma_2$,
and by \cref{prelim:subclasses::com:upper-sets:prop} it is commutative
so that  $\restrictAutGamS{\mathcal A}{\Gamma_1}{S}$
is a commutative permutation automaton.

$\ref{prelim:subclasses:thm:solvable:restrict} \Rightarrow \ref{prelim:prop:solvable-SCC:2}$:
This is clear.

$\ref{prelim:prop:solvable-SCC:2} \Rightarrow \ref{prelim:prop:solvable-SCC:3a}$:
Since a strongly connected component~$S$ is a block
by assumption, $\delta(S, a) \cap S \ne \emptyset$
and  $\delta(S, b) \cap S \ne \emptyset$
imply $\delta(S, a) = \delta(S, b) = S$ and then
$\delta(q, ab) = \delta(q, ba)$ for all $q \in S$.

$\ref{prelim:prop:solvable-SCC:3a} \Rightarrow \ref{prelim:prop:solvable-SCC:3}$:
Let $S \subseteq Q$ be strongly connected, $q \in \delta(S, a) \cap S$,
and $p \in \delta(S, b) \cap S$. Set $T = [q]_{\simeq}$.
Then $S \subseteq T$, $q \in \delta(T, a) \cap T$, $p \in \delta(T, b) \cap T$,
and so $\delta(r, ab) = \delta(r, ba)$ for all $r \in T$, in particular
for those states from~$S$.

$\ref{prelim:prop:solvable-SCC:3} \Rightarrow \ref{prelim:prop:solvable-SCC:4}$:
Let~$S$, $q \in S$ and $a,b \in \Sigma$ as in the statement.
Since $\delta(q, ab) \in S$ and $q \in S$, the set $T = S \cup \{\delta(q, a)\}$ is also strongly connected.
we have $\delta(q, a) \in T$ and $\delta(\delta(q, a), b) \in T$
so that, by Item~\ref{prelim:prop:solvable-SCC:3},
we find $\delta(q, ab) = \delta(q, ba)$.

$\ref{prelim:prop:solvable-SCC:4} \Rightarrow \ref{prelim:prop:solvable-SCC:3b}$:
Write $u = u_1 \cdots u_n$ with $u_i \in \Sigma$ for $i \in \{1,\ldots,n\}$.
Let $i \in \{1,\ldots,n\}$.
Since $\delta(q, u) = q$, we have 
\[
q \eqreach \delta(q, u_1\cdots u_{i-1}) \eqreach 
\delta(q, u_1\cdots u_{i-1} u_i) \eqreach \delta(q, u_1\cdots u_{i-1} u_i u_{i+1}).
\]
So, the set $S = \{q, \delta(q, u_1\cdots u_{i-1}), \delta(q, u_1\cdots u_{i-1} u_i),
\delta(q, u_1\cdots u_{i-1} u_i u_{i+1}), q\}$
is strongly connected. By assumption,
we get $\delta(p, u_i u_{i+1}) = \delta(p, u_{i+1}u_i)$
for all $p \in S$, in particular
we have
\[\delta(q, u_1\cdots u_{i-1} u_i u_{i+1})
= \delta(q, u_1\cdots u_{i-1} u_{i+1} u_i).
\]
Since every permutation of the letters $u_1, \ldots, u_n$
can be realized by swapping adjacent letters, we obtain the claim.

\ref{prelim:prop:solvable-SCC:3b} $\Rightarrow$ \ref{prelim:prop:solvable-SCC:3c}: Let $q \in Q$ and $u,v,w \in \Sigma^*$
be such that $\delta(q, u) = q$ and $u = vw$. Write $v = v_1 \cdots v_m$
for $v_1, \ldots, v_m \in \Sigma$ and 
consider a permutation $\pi \colon \{1,\ldots, m\} \to \{1,\ldots,m\}$.
Since $\delta(q, vw) = q$ we have $ \delta(q, v) = \delta(q, vwv)$
and $\delta(q, v_{\pi(1)} \cdots v_{\pi(m)})
= \delta(q, vwv_{\pi(1)} \cdots v_{\pi(m)})$.	
By assumption, the former equality gives
$\delta(q, vwv_{\pi(1)} \cdots v_{\pi(m)}) = \delta(q, v)$.
Hence, we find 
$\delta(q, v) = \delta(q, v_{\pi(1)}\cdots v_{\pi(m)})$.
This gives the claim.

$\ref{prelim:prop:solvable-SCC:3c} \Rightarrow \ref{prelim:classes:solvable:thm:char:solvable}$:
 Let $p \in Q$ and $a,b \in \Sigma$
 such that $\delta(q, ab) \lereach q$.
 Then $\delta(q, abu) = q$ for some $u \in \Sigma^*$
 and, by assumption,
 we have $\delta(q, ab) = \delta(q, ba)$.
\end{proof}

Our first corollary extends \cref{prelim:classes:comm:cor:collections-of-facts:7}.

\begin{corollary} \label{prelim:prop:strongly-trap-solvable}
	A strongly trap-connected (resp.~strongly connected) solvable
	automaton $\semiautomat{A}{Q}{\Sigma}{\delta}$ 
	is a commutative permutation-reset (resp.~commutative permutation) automaton.
\end{corollary}
\begin{proof}
	By \crefitem{prelim:automata:prop:char-strongly-trap}{prelim:automata:prop:char-strongly-trap:1b}
	we have $Q = [q]_{\eqreach} \cup \{t\}$ for a non-maximal $q \in Q$ and a trap state $t \in Q$ 
	(resp.~$Q = [q]_{\eqreach}$ for some $q \in Q$).
	Set $S = \upperset{\{q\}}$ and $T = S \setminus [q]_{\eqreach}$.
	Since $Q = S$ and $T = \{t\}$ (resp.~$T = \emptyset$)
	from 
	\crefitem{prelim:classes:solvable:thm}{prelim:classes:solvable:thm:char:SCC-quotients-comm-perm} 
	and \cref{prelim:Rees:S-at-most-one}
	we get
	that $\mathcal A$ is a commutative permutation-reset (resp.~commutative permutation)
	automaton.
\end{proof}

\begin{corollary} \label{cor:solvable-perm-SCC}
	Let $\semiautomat{A}{Q}{\Sigma}{\delta}$ be solvable
	and  $S \subseteq Q$ a non-empty union of sink components (as, for example, the set of maximal states).
	Then $\mathcal A_{|S}$ is a commutative permutation automaton.
\end{corollary}
\begin{proof}
	For each sink component $T \subseteq S$, $\mathcal A_{|T}$
	is solvable and strongly connected,
	and so a commutative permutation automaton
	by \cref{prelim:prop:strongly-trap-solvable}, which gives the claim. 
\end{proof}

\begin{corollary} Let $\semiautomat{A}{Q}{\Sigma}{\delta}$
	and $S \subseteq Q$ be an upper set.
	Then $\mathcal A / {\sim_S}$ is solvable
	if and only if for each strongly connected component
	$T \subseteq Q$ 
	we have $\delta(q, ab) = \delta(q, ba)$
	for all $a,b \in \Sigma$
	and $q \in T \setminus S$ with $\delta(q, ab) \in T$.
\end{corollary}
\begin{proof}
	This follows directly from \crefitem{prelim:classes:solvable:thm}{prelim:prop:solvable-SCC:4},
	the definition of the Rees quotient automaton,
	and the fact that a strongly connected set is either completely
	outside of~$S$ or completely inside of~$S$.	
\end{proof}

\begin{proposition}	\label{prelim:classes:solv:prop:char-seq}
Let $\semiautomat{A}{Q}{\Sigma}{\delta}$.
Then the following are equivalent.
\begin{enumerate}
	\item \label{prelim:thm:solvable:1} $\mathcal A$ is solvable. 
	
	\item \label{prelim:thm:solvable:2} There exists a sequence of upper sets
	$
	S_1 \subsetneq \cdots \subsetneq S_m = Q
	$
	where~$S_1$ equals the union of all sink components,
	$\mathcal A_{|S_1}$ is a commutative permutation automaton, and, for $i \in \{2,\ldots,m\}$,~$( \mathcal A_{|S_i} )/ {\sim_{S_{i-1}}}$ is
	a strongly trap-connected commutative permutation-reset automaton and $S_{i} \setminus S_{i-1}$ are all the strongly connected components that are not sink components.
		
	\item \label{prelim:classes:solv:prop:char-seq:3} There exists a sequence of upper sets
	$
	\emptyset = S_0 \subsetneq S_1 \subsetneq \cdots \subsetneq S_m = Q
	$
	such that~$( \mathcal A_{|S_i} )/ {\sim_{S_{i-1}}}$
	is a commutative permutation-reset automaton
	and $S_{i} \setminus S_{i-1}$ are all the strongly connected components
	for $i \in \{1,\ldots,m\}$. 
	
	\item \label{prelim:thm:solvable:solv-sequence} There exists a sequence of upper sets
	$
	\emptyset = S_0 \subsetneq S_1 \subsetneq \cdots \subsetneq S_m = Q
	$
	such that~$( \mathcal A_{|S_i} )/ {\sim_{S_{i-1}}}$ 
	is a solvable automaton for each $i \in \{1,\ldots,m\}$.   
\end{enumerate}
\end{proposition}
\begin{proof}
\ref{prelim:thm:solvable:1} $\Rightarrow$ \ref{prelim:thm:solvable:2}:
Consider a sequence of upper sets
as in \crefitem{prelim:general:order:cor:seq-upper-sets}{prelim:general:order:cor:seq-upper-sets:2}.
By \cref{cor:solvable-perm-SCC}, $\mathcal A_{|S_0}$ is a commutative permutation automaton.
Next, each set $S_i$ for $i \in \{1,\ldots, m\}$
has the form $S_i = [q]_{\simeq} \cup S_{i-1}$
for some $q \in S_i \setminus S_{i-1}$, where $q$
is not maximal since~$S_1$ equals the set of maximal states.
By \crefitem{prelim:classes:solvable:thm}{prelim:classes:solvable:thm:char:SCC-quotients-comm-perm},
$\mathcal ( \mathcal A_{|S_i} ) / {\sim_{S_{i-1}}}$ is a commutative permutation-reset automaton.\footnote{Alternatively, 
since the quotient automata
$\mathcal ( \mathcal A_{|S_i} ) / {\sim_{S_{i-1}}}$ for $i \in \{2,\ldots,m\}$
are strongly trap-connected by \cref{prelim:automata:prop:char-strongly-trap}
and solvable by \cref{prelim:classes:solvable:cor:subhom},
from \cref{prelim:prop:strongly-trap-solvable} we obtain 
that they are commutative permutation-reset automata.}
Since~$q$ is not a maximal state, $[q]_{\simeq}$
is not a sink component
and so $\mathcal ( \mathcal A_{|S_i} ) / {\sim_{S_{i-1}}}$
is strongly trap-connected by \cref{prelim:automata:prop:char-strongly-trap}.

\ref{prelim:thm:solvable:2} $\Rightarrow$ \ref{prelim:classes:solv:prop:char-seq:3}:
Consider a sequence $S_1 \subsetneq \cdots \subsetneq S_m = Q$
as in Item~\ref{prelim:thm:solvable:2}.
Applying \crefitem{prelim:general:order:cor:seq-upper-sets}{prelim:general:order:cor:seq-upper-sets:1} 
to $\mathcal A_{|S_1}$
there exists a sequence $\emptyset = T_0 \subsetneq T_1 \subsetneq \cdots \subsetneq T_l$
such that $T_i \setminus T_{i-1}$ for $i \in \{1,\ldots,l-1\}$
are all the sink components.
Since~$\mathcal A_{|S_1}$ is a commutative permutation automaton,
each $\mathcal A_{|T_i}$ is a commutative permutation automaton.
Then each $( \mathcal A_{|T_i} ) / {\sim_{T_{i-1}}}$ is a commutative permutation automaton.
Now, $T_0 \subsetneq T_1 \subsetneq \cdots \subsetneq T_{l} \subsetneq S_2 \subsetneq \cdots \subsetneq S_m$
yields a sequence as claimed.

\ref{prelim:classes:solv:prop:char-seq:3} $\Rightarrow$ \ref{prelim:thm:solvable:solv-sequence}:
As a commutative automaton is solvable, we obtain the claim.

\ref{prelim:thm:solvable:solv-sequence} $\Rightarrow$ \ref{prelim:thm:solvable:1}: Let $q \in Q$
and $a,b \in \Sigma$ be such that $\delta(q, ab) \lereach q$.
Since, by \cref{prelim:prop:reachability-finest}, $\delta(q,ab)$ and~$q$
are contained in the same strongly connected component, by \cref{prelim:general:order:subset-in-equiv},
we have $\{ \delta(q, ab), q \} \subseteq S_i \setminus S_{i-1}$ for some $i \in \{1,\ldots,m\}$.
Since $(\mathcal A_{|S_i} ) / {\sim_{S_i}}$ is solvable,
by \cref{prelim:rem:Rees} we have 
$\{ \delta_{|S_i}(q, ab) \} = \delta_{\sim_{S_i}}([q]_{\sim_{S_i}}, ab) = 
\delta_{\sim_{S_i}}([q]_{\sim_{S_i}}, ba) = \{ \delta_{|S_i}(q, ba) \}$
where $[q]_{\sim_{S_i}} = \{q\}$,
and so $\delta(q, ab) = \delta(q, ba)$.
Hence, $\mathcal A$ is solvable.
%
%$\ref{prelim:subclasses:thm:solvable:sequence} 
%\Rightarrow 
%\ref{prelim:subclasses:thm:solvable:solvable}$:  Let
%$
%\emptyset = S_0 \subsetneq S_1 \subsetneq \ldots \subsetneq S_m = Q
%$
%be a sequence of upper sets such that~$( \mathcal A_{|S_i} )/ {\sim_{S_{i-1}}}$
%is a commutative automaton for each $i \in \{1,\ldots,m\}$.
%We can assume $Q \ne \emptyset$, so that $m \ge 1$.
%For $m = 1$, $\mathcal A$ is isomorphic to~$( \mathcal A_{|S_1} )/ {\sim_{S_{0}}}$,
%which is commutative and so solvable.
%Otherwise,  $S_0 \subsetneq S_1 \subsetneq \ldots \subsetneq S_{m-1}$
%forms a sequence in $\mathcal A_{|S_{m-1}}$
%where, by \cref{prelim:prop:automata:isos:item:subaut,prelim:prop:automata:isos:item:corsub},
%$( ( \mathcal A_{|S_{m-1}} )_{|S_i} ) / {\sim_{S_{i-1}}} \cong ( \mathcal A_{|S_i} ) / {\sim_{S_{i-1}}}$
%are commutative for $i \in \{1,\ldots,m-1\}$.
%Inductively, we find that $\mathcal A_{|S_{m-1}}$
%is solvable.  As $\mathcal A / {\sim_{S_{m-1}}}$
%is commutative, $\mathcal A$ is solvable.
\end{proof}

\Cref{prelim:classes:fig:exsolvable} gives a commutative automaton
with a sequence as in \crefitem{prelim:classes:solv:prop:char-seq}{prelim:thm:solvable:2}.

\begin{figure}[ht]
	\centering
	\begin{subfigure}{0.3\textwidth}
		\centering
		\begin{tikzpicture}[node distance = 2cm, on grid, auto, ->]
			\node (q1) [state] {$q_1$};
			\node (q2) [state, above of = q1] {$q_2$};
			
			\path[->] (q1) edge node {$a$} (q2)
			(q2) edge [loop above] node {$a$} (q2);
			%	(q1) edge [loop left] node {$b$} (q1);
		\end{tikzpicture}
		\caption{An automaton $\mathcal A$.}
	\end{subfigure}%
	\begin{subfigure}{0.3\textwidth}
		\centering
		\begin{tikzpicture}[node distance = 2cm, on grid, auto, ->]
			\tikzstyle{state-subset}=[]
			
			\node (11) [state-subset] {$(q_1, q_1)$};
			\node (12) [state-subset, right of = 11] {$(q_1, q_2)$};
			\node (21) [state-subset, above of = 11] {$(q_2, q_1)$};
			\node (22) [state-subset, above of = 12] {$(q_2, q_2)$};
			
			\path[->] (11) edge node {$a$} (22)
			(12) edge node {$a$} (22)
			(21) edge node {$a$} (22)
			(22) edge [loop above] node {$a$} (22);
			%(11) edge [loop left] node {$b$} (11)
			%(12) edge [loop right] node {$b$} (12)
			%(21) edge [loop left] node {$b$} (21);		 
		\end{tikzpicture}
		\caption{The automaton $\mathcal A \times \mathcal A$.}
	\end{subfigure}%
	\begin{subfigure}{0.35\textwidth}
		\centering
		\begin{align*}
			S_0 & = \{ (q_2, q_2) \}, \\
			S_1 & = \{ (q_2, q_2), (q_1, q_2)\}, \\
			S_2 & = \{ (q_2, q_2), (q_1, q_2), (q_1, q_1) \}, \\
			S_4 & = \{ q_1, q_2\} \times \{ q_1, q_2\}.
		\end{align*}		
		\caption{Sequence of upper sets.}
	\end{subfigure}	
	\caption[Solvable automaton with sequence of upper sets]{A sequence of subsets $S_0 \subseteq S_1 \subseteq S_2 \subseteq S_3$
		as stated in \crefitem{prelim:classes:solv:prop:char-seq}{prelim:thm:solvable:2}
		for a unary (and so solvable) automaton of the form $\mathcal A \times \mathcal A$.}
	\label{prelim:classes:fig:exsolvable}
\end{figure} 

\begin{proposition} \label{prelim:classes:solvable:prop:char-inductive}
Let $\semiautomat{A}{Q}{\Sigma}{\delta}$. Then the following are equivalent.
\begin{enumerate}
\item \label{prelim:classes:solvable:prop:char-inductive:1} $\mathcal A$ is solvable.	
	
\item \label{prelim:classes:solvable:prop:char-inductive:2} $\mathcal A$ is a commutative permutation-reset automaton or there exists a proper upper set $S \subseteq Q$
with $|S| \ge 2$ such that $\mathcal A_{|S}$ is solvable and $\mathcal A / {\sim_S}$ is a commutative
and strongly trap-connected permutation-reset automaton.

\item \label{prelim:classes:solvable:prop:char-inductive:3} $\mathcal A$ is a commutative permutation-reset automaton or there exists a proper upper set $S \subseteq Q$
with $|S| \ge 2$ such that $\mathcal A_{|S}$ is a commutative permutation-reset automaton
and $\mathcal A / {\sim_S}$ is solvable.
	
\item \label{prelim:classes:solvable:prop:char-inductive:4} $\mathcal A$ is commutative or there exists a proper upper set $S \subseteq Q$ with $|S| \ge 2$
such that $\mathcal A_{|S}$ and $\mathcal A / {\sim_S}$ are solvable. 
	
\item \label{prelim:classes:solvable:prop:char-inductive:5} There exists an upper set $S \subseteq Q$
such that $\mathcal A_{|S}$ and $\mathcal A / {\sim_S}$ are solvable. 
\end{enumerate}
\end{proposition}
\begin{proof} \ref{prelim:classes:solvable:prop:char-inductive:1} $\Rightarrow$ \ref{prelim:classes:solvable:prop:char-inductive:2} \& \ref{prelim:classes:solvable:prop:char-inductive:1} $\Rightarrow$ \ref{prelim:classes:solvable:prop:char-inductive:3}: Consider a sequence 
$
S_1 \subsetneq \cdots \subsetneq S_m = Q
$
as in \crefitem{prelim:classes:solv:prop:char-seq}{prelim:thm:solvable:2}. For $m = 1$, $Q$ equals the set
of maximal states and~$\mathcal A$ is a commutative permutation automaton
and so Items~\ref{prelim:classes:solvable:prop:char-inductive:2} and~\ref{prelim:classes:solvable:prop:char-inductive:3} hold.
Next, suppose $m \ge 2$.
Then $\mathcal A / {\sim_{S_{m-1}}}$ is a 
commutative and strongly trap-connected permutation-reset automaton.
For $|S_{m-1}| = 1$, we have $\mathcal A \cong \mathcal A / {\sim_{S_{m-1}}}$
and this shows Item~\ref{prelim:classes:solvable:prop:char-inductive:2}.
Otherwise, since $\mathcal A_{|S_{m-1}}$ is solvable by \cref{prelim:classes:solvable:prop:subaut},
Item~\ref{prelim:classes:solvable:prop:char-inductive:2} holds
with $S = S_{m-1}$.

Next, if $|S_1| \ge 2$, then Item~\ref{prelim:classes:solvable:prop:char-inductive:3} 
is established since $\mathcal A_{|S_1}$ is a commutative permutation
automaton and $\mathcal A / {\sim_{S_1}}$
is solvable by \cref{prelim:classes:solvable:cor:subhom}. For $|S_1| = 1$,
we have $\mathcal A_{|S_2} \cong ( \mathcal A_{|S_2} ) / {\sim_{S_1}}$,
and the latter automaton is a commutative permutation-reset automaton
by choice of the sequence. 
Thus for $m = 2$, $\mathcal A$ is a commutative permutation-reset automaton
and Item~\ref{prelim:classes:solvable:prop:char-inductive:3} holds
or Item~\ref{prelim:classes:solvable:prop:char-inductive:3} holds with the proper upper set of size at least two $S = S_2$.

\ref{prelim:classes:solvable:prop:char-inductive:2} $\Rightarrow$ \ref{prelim:classes:solvable:prop:char-inductive:4} \& \ref{prelim:classes:solvable:prop:char-inductive:3} $\Rightarrow$ \ref{prelim:classes:solvable:prop:char-inductive:4} \& \ref{prelim:classes:solvable:prop:char-inductive:4} $\Rightarrow$ \ref{prelim:classes:solvable:prop:char-inductive:5}: These implications are clear.
		
\ref{prelim:classes:solvable:prop:char-inductive:5} $\Rightarrow$ \ref{prelim:classes:solvable:prop:char-inductive:1}: Let $S \subseteq Q$
be an upper set such that $\mathcal A_{|S}$ and $\mathcal A / {\sim_S}$ are solvable. We can assume $0 < |S| < |Q|$. 
Then $\emptyset \subsetneq S \subsetneq Q$ 
gives a sequence as in \crefitem{prelim:classes:solv:prop:char-seq}{prelim:thm:solvable:solv-sequence}.
So $\mathcal A$ is solvable.
\end{proof}

\Crefitem{prelim:classes:solvable:prop:char-inductive}{prelim:classes:solvable:prop:char-inductive:2} 
is similar to \cref{prelim:automata:cor:composition-series-maximal-states},
but gives more information on automata-theoretical properties.

\begin{remark}
\Crefitem{prelim:classes:solvable:prop:char-inductive}{prelim:classes:solvable:prop:char-inductive:4}
can be used to give an alternative inductive 
proof of the statement from \cref{prelim:subclasses:prop:solvable-product} that $\mathcal A_1 \times \mathcal A_2$ is solvable
for solvable~$\mathcal A_1, \mathcal A_2$.
Since, if $\mathcal A_1$ and $\mathcal A_2$
are commutative, then it is readily seen that $\mathcal A_1 \times \mathcal A_2$
is commutative.
Next, if there exists, without loss of generality,
an upper set $S_1 \subseteq Q_1$
with $2 \le |S_1| < |Q_1|$
such that $( \mathcal A_1 )_{|S_1}$
and $\mathcal A_1 / {\sim_{S_1}}$ are solvable, then the set $T = S_1 \times Q_2$
is an upper set in $\mathcal A_1 \times \mathcal A_2$
and by \cref{prelim:automata:prod:eqn:sub,prelim:automata:prod:eqn:quot} we have
\[
( \mathcal A_1 \times \mathcal A_2 )_{|T} \cong (\mathcal A_1)_{|S_1} \times \mathcal A_2 \quad 
\text{ and }
\quad
( \mathcal A_1 \times \mathcal A_2 ) / {\sim_T} \cong \mathcal A_1 / {\sim_{S_1}}.
\]
Inductively, $(\mathcal A_1)_{|S_1} \times \mathcal A_2$
is solvable. So we obtain solvability of $\mathcal A_1 \times \mathcal A_2$.
\end{remark}

\begin{remark}
	\Crefitem{prelim:classes:solvable:prop:char-inductive}{prelim:classes:solvable:prop:char-inductive:4}
	can be used to give an alternative proof
	for \cref{prelim:classes:solvable:prop:subaut}
	using the various isomorphisms established in \cref{prelim:sec:rees}
	as we shown next with the same notation as in \cref{prelim:classes:solvable:prop:subaut}.	
	First, commutativity of~$\mathcal A$
	transfers to $\mathcal A_{|S}$.
	Next, suppose
	there exists an upper set $R \subseteq Q$
	with $2 \le |R| < |Q|$
	such that~$\mathcal A_{|R}$
	and $\mathcal A / {\sim_R}$ are solvable. 
	By \cref{prelim:prop:automata:isos:item:subaut} we have
	\[
	( \mathcal A_{|S} )_{|S \cap R} \cong \mathcal A_{|S \cap R} \cong ( \mathcal A_{|R} )_{|S \cap R},
	\]
	and with the canonical projection map  $\pi_{\sim_R} \colon Q \to Q / {\sim_R}$ introduced in \cref{prelim:sec:relations}, 
	by \cref{prelim:eqn:Rees:ScapT}, %,prelim:eqn:Rees:ScapT}, 
\[
%(\mathcal A_{|T}) / {\sim_{T\cap R}} \cong ( \mathcal A / {\sim_{T\cap R}})_{|\pi_{\sim_{T\cap R}}(T)} \cong (\mathcal A / {\sim_R})_{|\pi_{\sim_R}(T)}.
(\mathcal A_{|S}) / {\sim_{S\cap R}} \cong (\mathcal A / {\sim_R})_{|\pi_{\sim_R}(S)}.
\]
Inductively, $( \mathcal A_{|R} )_{|S \cap R}$ 
and $(\mathcal A / {\sim_R})_{|\pi_{\sim_R}(S)}$ are solvable.
So, by \crefitem{prelim:classes:solvable:prop:char-inductive}{prelim:classes:solvable:prop:char-inductive:4}
or by \cref{prelim:Rees:S-at-most-one}, $\mathcal A_{|S}$ is solvable.	
\end{remark}

\begin{remark}
	\Crefitem{prelim:classes:solvable:prop:char-inductive}{prelim:classes:solvable:prop:char-inductive:4}
	can be used to give an alternative proof
	for \cref{prelim:classes:solvable:prop:morph}
	using the various isomorphisms established in \cref{prelim:sec:rees}
	as we shown next with the same notation as in \cref{prelim:classes:solvable:prop:morph}.
	First, if $\mathcal A$ is commutative,
	then clearly $\mathcal B$ is commutative.
	Now, suppose
	there exists an upper set $R \subseteq Q$
	with $2 \le |R| < |Q|$
	such that~$\mathcal A_{|R}$
	and $\mathcal A / {\sim_R}$ are solvable.
	The set $\varphi(R)$ is an upper set in $\mathcal B$ by \cref{prelim:automata:morpisms:reachability:prop:images}.
	Now,~$\varphi_{|R}$ considered as a surjective function from $R$ onto $\varphi(R)$
	is a surjective morphism from~$\mathcal A_{|R}$ onto~$\mathcal B_{|\varphi(R)}$.
	%By \crefitem{prelim:thm:homtheorem}{prelim:thm:homtheorem:4}
	By \cref{prelim:prop:Rees-homsfactorthrough}
	there exists a surjective morphism
	from~$\mathcal A / {\sim_R}$ onto~$( \mathcal B / {\sim_{\varphi(R)}} )_{|\pi_{\sim_{\varphi(R)}}(\varphi(Q))}$,
	where $( \mathcal B / {\sim_{\varphi(R)}} )_{|\pi_{\sim_{\varphi(R)}}(\varphi(Q))} = \mathcal B / {\sim_{\varphi(R)}}$
	as $\varphi(Q) = S$.
	By induction hypothesis, 
	as~$\mathcal A_{|R}$ and~$\mathcal A / {\sim_R}$ are solvable,
	$\mathcal B_{|\varphi(R)}$
	and $\mathcal B / {\sim_{\varphi(R)}}$ are solvable.
	So, by \crefitem{prelim:classes:solvable:prop:char-inductive}{prelim:classes:solvable:prop:char-inductive:4},
	$\mathcal B$ is solvable.
\end{remark}

\subsection{Synchronization and Subset Mapping Properties}
\label{prelim:classes:sec:solvable:sync}

First, in \cref{constr_sync:cor:solvable_sync_iff_reach_trap},
we state that in solvable automata, a synchronizing state is a trap state (and hence unique).
\Cref{prelim:thm:solvrank_r} states the bound~$n - r$
for shortest words of rank at most~$r$ in solvable automata that admit a word of rank~$r$,
thus extending Rystsov's bound for $r = 1$~\cite[Cor.~4.2]{Rystsov97}.
Proposition~\ref{prop:S-T-bound}
gives words of a specific form to map
one subset onto another subset
in solvable automata. This result is rather technical, and we need it
in \cref{compl:sec:solv} (\cref{thm:solvable}) to prove that \ST\ 
is in \NP\ for solvable automata.

\begin{proposition}
	\label{constr_sync:cor:solvable_sync_iff_reach_trap} \label{constr_sync:prop:solvable_sync_iff_reach_trap}
	A synchronizing state in a solvable automaton
	is a trap state.	
\end{proposition}
\begin{proof}
	By \crefitem{prelim:classes:solvable:thm}{prelim:prop:solvable-SCC:2}
	in solvable automata the sink components are blocks.
	So, the claim follows from \cref{constr_sync:prop:sink-blocks-sync-trap}.  
\end{proof}

\begin{theorem}
	\label{prelim:thm:solvrank_r}
	Let $\semiautomat{A}{Q}{\Sigma}{\delta}$ be an $n$-state solvable
	automaton. If there exists a word of rank $r$ in $\mathcal A$,
	then there exists a word of length
	at most $n - r$ whose rank is at most $r$ 
	and these bounds are tight even for unary automata where~$r$
	is the minimal rank among all words. Furthermore, for $\Sigma = \{a_1, \ldots, a_k\}$,
	the word $(a_1 \cdots a_k)^{n - r}$
	has rank at most~$r$.
\end{theorem}
\begin{proof} 
	By \crefitem{prelim:classes:solvable:thm}{prelim:prop:solvable-SCC:2}
	the strongly connected components in $\mathcal A$
	are blocks.
	So, Theorem~\ref{prelim:thm:SCCblocks-rank_r}
	gives the upper bound $n - r$
	and that $(a_1 \cdots a_k)^{n - r}$ 
	has rank at most~$r$.
	Tightness as claimed follows from the automata
	given in the proof of \cref{prelim:prop:unary_bound}.
\end{proof}

\begin{proposition} \label{prop:S-T-bound}
	Let $\semiautomat{A}{Q}{\Sigma}{\delta}$ be solvable with $n$ states, $S \subseteq Q$,
	and $w \in \Sigma^*$.
	Then for $\Sigma = \{a_1, \ldots, a_k\}$
	there exists $v \in \Sigma^*$ with $|v| < 2^n$
	and $\delta(S, v) = \delta(S, w)$
	of the form
	\[ 
	v = v_0 x_1 \cdots v_{r-1} x_r v_r
	\]%  a_1^{r_{0,1}} \cdots a_k^{r_{0,k}}
	for some $r \le n^2$
	where $x_1, \ldots, x_r \in \Sigma$ and $v_i = a_1^{r_{i,1}} \cdots a_k^{r_{i,k}}$
	for numbers $0 \le r_{i,j} < 2^n$
	with $i \in \{0,1,\ldots,r\}$ and $j \in \{1,\ldots,k\}$.
\end{proposition}
\begin{proof}
	We can assume $S \ne \emptyset$.
	By \crefitem{prelim:classes:solv:prop:char-seq}{prelim:thm:solvable:2} 
	there exist upper sets 
	\begin{equation}\label{eqn:S-T:sequence}
		S_1 \subsetneq \cdots \subsetneq S_m = Q
	\end{equation}
	where~$S_1$ equals the union of all sink components,
	$\mathcal A_{|S_1}$ is commutative, and $S_{i} \setminus S_{i-1}$ for $i \in \{2,\ldots,m\}$
	are all the strongly connected components that are not sink components.  
	Consider a word
	$u \in \Sigma^*$ of shortest possible
	length with $\delta(S, u) = \delta(S, w)$.
	Let $S = \{s_1, \ldots, s_l\}$ for $l = |S| \ge 1$
	and set $P = \pref{\{u\}}$.
	Next, we define a function $f \colon P \to \mathbb \{1,\ldots,m\}^l$ 
	by setting $f(u') = (i_1,\ldots,i_l)$  for $u' \in P$ 
	where~$i_j \in \{1,\ldots,m\}$ equals the minimal number with $\delta(s_j, u') \in S_{i_j}$
	for all $j \in \{1,\ldots,l\}$.
	%by setting for $u' \in \Sigma^*$ such that $u = u' u''$ for some $u'' \in \Sigma^*$
	%let $f(u') = (i_1,\ldots,i_l)$ 
	%be the vector where $i_j$ is minimal such that $\delta(s_j, u') \in S_{i_j}$
	%for all $j \in \{1,\ldots,l\}$.
	Now, there exists a unique factorization $u = u_0 x_1 u_1 \cdots u_{r-1} x_r u_r$
	with  $u_0, u_1, \ldots, u_r \in \Sigma^*$
	and $x_1, \ldots, x_r \in \Sigma$ for some $r \ge 0$
	such that $f(u_0 x_1 \cdots u_{t-1}) \ne f(u_0 x_1 \cdots u_{t-1}x_t)$
	for $t \in \{1,\ldots,r\}$ and
	\begin{align*}
		f(\varepsilon) & = f(u_0), \\
		f(u_0x_1)      & = f(u_0x_1 u_1), \\
		\vdots \\
		f(u_0 x_1 \cdots u_{r-1} x_r) & = f(u_0 x_1 \cdots u_{r-1} x_r u_r).
	\end{align*}
	Since the $S_i$ are upper sets,
	for $t \in \{1,\ldots,r\}$ with $f(u_0 x_1 \cdots u_{t-1}) = (i_1, \ldots, i_l)$ 
	and 
	\[ 
	f(u_0 x_1 \cdots u_{t-1} x_t) = (i_1', \ldots, i_l')
	\] 
	we have $i_j \ge i_j'$ for all $j \in \{1,\ldots,l\}$
	where $i_j > i_j'$ for at least one $j \in \{1,\ldots,l\}$. 
	So, $i_1' + \ldots + i_l' < i_1 + \ldots + i_l$.
	As these sums are between~$l$ and~$lm$, they can decrease at most $l(m-1)$ times 
	and we find $r \le l(m-1)$. As $l(m-1) \le n^2$ we get $r \le n^2$. 
	Now, for $j \in \{1,\ldots,l\}$ and $t \in \{0,1,\ldots,r\}$
	the states $\delta(s_j, u_0 x_1 \cdots u_{t-1} x_t)$ and 
	$\delta(s_j, u_0 x_1 \cdots u_{t-1} x_t u_t)$ are both contained in the strongly connected component
	$S_{i_j} \setminus S_{i_j - 1}$ for $i_j \ge 2$
	or both states are contained in $S_1$.\footnote{We set $u_0 x_1 \cdots u_{-1} x_0 = \varepsilon$.}
	By \crefitem{prelim:classes:solvable:thm}{prelim:prop:solvable-SCC:3b}
	and as $\mathcal A_{|S_1}$ is commutative 
	we can reorder the letters in~$u_t$ while still ending up at the same state
	and so we find
	\[ 
	\delta(s_j, u_0 x_1 \cdots u_{t-1} x_t u_t) = \delta(s_j, u_0 x_1 \cdots u_{t-1} x_t a_1^{|u_t|_{a_1}} \cdots a_k^{|u_t|_{a_k}}).
	\]
	Now, as we have chosen $u$ of minimal length,
	for distinct $u',u''\in P$ we have $\delta(S, u') \ne \delta(S, u'')$
	as otherwise, for, without loss of generality, $|u'| < |u''|$
	we have $\delta(S, u'u''') = \delta(S, u'' u''')$ 
	for $u''' \in \Sigma^*$ with $u = u'' u'''$
	where $|u' u'''| < |u|$, contradicting the minimal choice of~$u$.
	%(i.e., $u$ traces out a path in the power automaton that has the subsets of~$Q$ as states
	%and transitions for each letter given by the action of $\delta$ on subsets without repeating any sets) where these set are nonempty.
	So, each $u' \in P$ corresponds to a distinct nonempty subset of states $\delta(S, u')$
	and we conclude $|u| \le 2^n - 1$, which gives
	$|u_t|_{a_h} \le 2^n - 1$ for each $h \in \{1,\ldots,k\}$.
	Then for the word
	$v = v_0 x_1 \cdots v_{r-1} x_r v_r$
	where $v_i = a_1^{|u_i|_{a_1}} \cdots a_k^{|u_i|_{a_k}}$
	for $i \in \{0,1,\ldots,r\}$
	we have $\delta(S, u) = \delta(S, v)$
	and the claim is proved with~$v$.
\end{proof}

\subsection{Comparison with Original Definition}
\label{prelim:classes:sec:comparison}
	
\IN{$R_1$-solvable}
\IN{$R_2$-solvable}
Rystsov~\cite{Rys96} originally defined an automaton to be solvable, if
it is synchronizing and admits a sequence $S_1 \subsetneq \cdots \subsetneq S_m = Q$ 
of upper sets such that $S_1$ is a singleton set 
containing a trap state and $\mathcal A_{|S_i} / {\sim_{S_{i-1}}}$ for $i \in \{2,\ldots,m\}$ are commutative. 
This is equivalent to being synchronizing and solvable
since, with \cref{constr_sync:cor:solvable_sync_iff_reach_trap}, it follows that~$S_1$ is always a singleton 
for synchronizing and solvable automata.
We call a synchronizing and solvable automaton
also an~\emph{$R_1$-solvable} automaton in the following.
Furthermore, we say that~$\semiautomat{A}{Q}{\Sigma}{\delta}$
is \emph{$R_2$-solvable} if $(\mathcal A \times \mathcal A) / {\sim_D}$
is $R_1$-solvable where $D = \{(q,q) \mid q \in Q\}$. This was
another notion of solvability considered in~\cite{Rys96}.
Since synchronizability of $(\mathcal A \times \mathcal A) / {\sim_D}$
is equivalent to synchronizability of $\mathcal A$ by \cref{prelim:subclasses:sync:Aiffpairaut},
an $R_2$-solvable automaton is synchronizing.\footnote{In the original
	definition from~\cite[p.~277]{Rys96} 
	synchronizability of~$\mathcal A$ was explicitly included.}
	
\begin{proposition} \label{prelim:classes:solvable:prop:R1}
	Let $\semiautomat{A}{Q}{\Sigma}{\delta}$
	and $S \subseteq Q$ be the set of maximal states in~$\mathcal A$ w.r.t.~$\lereach$.
	Then $\mathcal A$ is solvable if and only if $\mathcal A / {\sim_S}$
	is $R_1$-solvable and $\mathcal A_{|S}$ is a commutative permutation automaton.
\end{proposition}
\begin{proof}
	First, suppose~$\mathcal A$ is solvable.
	Then, by \crefitem{prelim:classes:solv:prop:char-seq}{prelim:thm:solvable:2}, 
	$\mathcal A_{|S}$ is a commutative permutation automaton. 
	By \cref{prelim:classes:solvable:cor:subhom}, $\mathcal A / {\sim_S}$
	is solvable and, by \cref{prelim:prop:maximal-states-factor-sync}, 
	it is synchronizing; and so $R_1$-solvable.
	
	Conversely, suppose $\mathcal A_{|S}$ is a commutative permutation automaton
	and	$\mathcal A / {\sim_S}$ is $R_1$-solvable, and so also solvable.
	Thus, by \crefitem{prelim:classes:solvable:prop:char-inductive}{prelim:classes:solvable:prop:char-inductive:5}, 
	since also $\mathcal A_{|S}$ is solvable, $\mathcal A$ is solvable.	
\end{proof}

\begin{proposition}
	An automaton $\semiautomat{A}{Q}{\Sigma}{\delta}$ is $R_2$-solvable
	if and only if it is synchronizing and $(\mathcal A \times \mathcal A) / {\sim_D}$
	with $D = \{(q,q) \mid q \in Q\}$ is solvable.
\end{proposition}
\begin{proof}
    If~$\mathcal A$ is $R_2$-solvable, then $(\mathcal A \times \mathcal A) / {\sim_D}$
    is $R_1$-solvable, and so solvable and synchronizing.
    Synchronizability of $(\mathcal A \times \mathcal A) / {\sim_D}$
    implies synchronizability of $\mathcal A$ by \cref{prelim:subclasses:sync:Aiffpairaut}.
        
    Conversely, \cref{prelim:subclasses:sync:Aiffpairaut}
    implies synchronizability of $(\mathcal A \times \mathcal A) / {\sim_D}$
    and so it is $R_1$-solvable. Thus~$\mathcal A$ is $R_2$-solvable.
\end{proof}

\begin{proposition} \label{prelim:classes:solvable:prop:R1-implies-R2}
	Let $\semiautomat{A}{Q}{\Sigma}{\delta}$.
	If $\mathcal A$ is $R_1$-solvable,
	then $\mathcal A$ is $R_2$-solvable.
\end{proposition}
\begin{proof}
    Since~$\mathcal A$ is solvable, \cref{prelim:subclasses:cor:solv-pair} gives solvability
	of $(\mathcal A \times \mathcal A ) / {\sim_D}$.
	Then, synchronizability of $\mathcal A$ yields that it is $R_2$-solvable.
\end{proof}

By the following corollary, the class of $R_1$-solvable automata also equals
the intersection of the class of $R_2$-solvable with the class of solvable automata.

\begin{corollary}
	The automaton $\semiautomat{A}{Q}{\Sigma}{\delta}$
	is $R_1$-solvable if and only if it is $R_2$-solvable and solvable.
\end{corollary}
\begin{proof}
 An $R_1$-solvable automaton is solvable and, by \cref{prelim:classes:solvable:prop:R1-implies-R2},
 $R_2$-solvable.
 Conversely, since an $R_2$-solvable automaton is synchronizing,
 if it is solvable, it is also $R_1$-solvable.	
\end{proof}

\Cref{prelim:classes:fig:solvablenotR2solvable} 
depicts an $R_2$-solvable automaton that is not solvable, and so also not $R_1$-solvable.
By definition, every non-synchronizing solvable automaton
is neither $R_1$- nor $R_2$-solvable. 
Figure~\ref{fig:solvability-notions} shows the inclusions relations
between our notions of solvability that follow from the above considerations.

\begin{figure}[h]
	\centering
\begin{tikzpicture}	
	\begin{scope}[shift={(3cm,-5cm)}, fill opacity=0.5,	mytext/.style={text opacity=1}]
		
		%\draw[fill=red, draw = black] (0,0) circle (5);
		\draw[fill=gray, draw = black,name path=circle 1] (-2.5,0) ellipse (4.7 and 2);
		\draw[fill=gray,     draw = black,name path=circle 2] (2.5,0)  ellipse (3 and 2);
		
		\pgftransparencygroup
		\clip (-2.5,0) ellipse (4.7 and 2);
		\fill[gray!0.5] (2.5,0) ellipse (3 and 2);
		\filldraw[draw,fill={gray!0.5},name intersections={of=circle 1 and circle 2}] (intersection-1) .. controls +(-6,1) and +(-6,-1) ..(intersection-2);
		\endpgftransparencygroup
		
		%\node[mytext] at (0,4) (A) {A};
		\node[mytext] at (-5.5,0) (B) {synchronizing};
		\node[mytext] at (3.8,0) (C) {solvable};
		\node[mytext] at (0.8,0) (D) {$R_1$-solvable};
		\node[mytext] at (-2,0) (E) {$R_2$-solvable};
	\end{scope}
\end{tikzpicture}
	\caption[Inclusion relations for solvability notions]{Inclusion relations between the different solvability notions.}
	\label{fig:solvability-notions}
\end{figure}

\begin{remark}
	Yet another, more general notion,
    is obtained by considering
	those automata where $(\mathcal A \times \mathcal A) / {\sim_D}$
	is solvable.
	This class contains the $R_2$-solvable automata
	and, by \cref{prelim:subclasses:cor:solv-pair},
	all solvable automata.
	By adding a trap state to the automaton~$\mathcal A$
	from \cref{prelim:classes:fig:solvablenotR2solvable},
	we obtain a non-solvable automaton
	such that $(\mathcal A \times \mathcal A) / {\sim_D}$
	is solvable but not synchronizing.
	Thus this class properly
	contains the $R_2$-solvable automata.	
\end{remark}